\chapter*{پیش‌گفتار}
\addcontentsline{toc}{chapter}{پیش‌گفتار}
\markboth{پیش‌گفتار}{پیش‌گفتار}

این کتاب برای شماست؛ برای کسی که تا امروز حتی یک خط برنامه ننوشته است و
می‌خواهد از صفر شروع کند. هیچ پیش‌نیازی لازم نیست: نه ریاضیِ پیشرفته، نه
زبانِ انگلیسی، نه آشناییِ قبلی با رایانه فراتر از روشن و خاموش کردنِ آن.

\section*{این کتاب دربارهٔ چیست؟}

این کتاب دربارهٔ \emph{خودِ برنامه‌نویسی} است: اینکه چگونه یک مسئله را به
گام‌های کوچک بشکنیم (به این کار \textbf{تفکرِ الگوریتمی} می‌گویند)، و چگونه آن
گام‌ها را با دستورهای سادهٔ یک زبانِ برنامه‌نویسی بنویسیم تا رایانه اجرایشان
کند. در این کتاب فقط با سنگِ بنای برنامه‌نویسی سروکار داریم:

\begin{itemize}
  \item الگوریتم و حلِ مسئله؛
  \item متغیرها و انواعِ داده؛
  \item محاسبه با عملگرها؛
  \item تصمیم‌گیری با شرط‌ها؛
  \item تکرار با حلقه‌ها؛
  \item تقسیمِ برنامه به توابع؛
  \item نگهداریِ چند مقدار در آرایه‌ها؛
  \item و در پایان، چند الگوریتمِ کلاسیک که هر برنامه‌نویسی باید بشناسد.
\end{itemize}

همین و بس. خبری از موضوع‌های پیشرفته، کتابخانه‌های جانبی، طراحیِ وب یا
رابطِ گرافیکی نیست؛ آن‌ها همه بعد از این کتاب می‌آیند. هدفِ ما این است که
وقتی صفحهٔ آخر را بستید، \emph{برنامه‌نویس فکر کنید}: مسئله‌ای ببینید و
ناخودآگاه آن را به گام‌های قابلِ اجرا بشکنید.

\section*{چرا زبانِ سلام؟}

برای نوشتنِ برنامه به یک زبانِ برنامه‌نویسی نیاز داریم. در این کتاب از
\textbf{سلام} استفاده می‌کنیم، زبانی که می‌گذارد دستورها را به فارسی
بنویسید. وقتی مانعِ زبانِ بیگانه برداشته شود، تمامِ توانِ شما صرفِ یادگیریِ
\emph{مفهوم}‌ها می‌شود، نه حفظِ واژه‌های ناآشنا. آنچه اینجا می‌آموزید
(متغیر، شرط، حلقه، تابع، الگوریتم) در همهٔ زبان‌های برنامه‌نویسیِ دنیا همین
است؛ فقط لباسِ واژه‌ها فرق می‌کند. پس این کتاب شما را برنامه‌نویس می‌کند،
نه فقط «برنامه‌نویسِ سلام».

\section*{چگونه این کتاب را بخوانیم؟}

فصل‌ها را به ترتیب بخوانید؛ هر فصل بر فصلِ پیش استوار است. در سراسرِ کتاب
با این نشانه‌ها روبه‌رو می‌شوید:

\begin{note}
نکته‌ای برای روشن‌تر شدنِ مطلب یا یادآوریِ چیزی مهم.
\end{note}

\begin{tip}
ترفند یا روشی که کارِ شما را آسان‌تر می‌کند.
\end{tip}

\begin{warn}
هشدار دربارهٔ اشتباهی رایج که تازه‌کارها معمولاً گرفتارش می‌شوند.
\end{warn}

\begin{deep}
توضیحی ژرف‌تر دربارهٔ اینکه چیزی \emph{چرا} این‌گونه است. خواندنش اختیاری
است؛ می‌توانید از آن بگذرید و بعدها بازگردید.
\end{deep}

در پایانِ هر فصل چند \textbf{تمرین} آمده است و پاسخِ همهٔ آن‌ها در پیوستِ
انتهای کتاب. اما یک قولِ مردانه بدهید: پیش از دیدنِ پاسخ، خودتان تلاش کنید.
برنامه‌نویسی مانندِ شنا و دوچرخه‌سواری است؛ با \emph{انجام دادن} یاد گرفته
می‌شود، نه با تماشا کردن.

\vspace{0.5cm}
\begin{center}
\emph{آماده‌اید؟ بیایید یاد بگیریم چگونه با رایانه حرف بزنیم.}
\end{center}
\clearpage
